% 05_capi_s2s_deliverables

\section{Income Imputation (Survey-to-Survey)}

The L2Phl 2025 baseline does not field a full consumption module, so
household welfare cannot be measured directly. Instead it is predicted
through survey-to-survey (S2S) imputation from the Family Income and
Expenditure Survey--Labor Force Survey (FIES--LFS) 2023, which observes
welfare directly. Poverty is then a downstream indicator: predicted
welfare is compared against regional poverty lines. The method and
diagnostics summarized here follow the technical note prepared by
Sandra Segovia Juarez.

\subsection{Method}

The imputation proceeds in three steps.

\begin{procbox}{S2S imputation}
\textbf{1. Estimate the welfare model.} A weighted OLS model is fit in
FIES--LFS 2023 with log per-capita income as the dependent variable and
predictors available in both surveys: demographic composition, household
size, assets, housing and sanitation conditions, tenure, urban location,
and region fixed effects. The model uses 162{,}687 observations, with an
R-squared of 0.5108 and a root MSE of 0.4798.

\textbf{2. Apply the coefficients.} The estimated coefficients are
applied to the harmonized L2Phl variables to predict welfare for each
household.

\textbf{3. Add stochastic residuals.} Resampled source-survey residuals
are added to each prediction, so the imputed welfare distribution is not
artificially compressed.
\end{procbox}

Harmonizing the predictor definitions across the two surveys is the
critical step, and the most common source of error.

\subsection{GDP adjustment}

Because the model captures 2023 relationships, predicted welfare is
scaled up by cumulative real GDP growth over 2024--2025:
$(1 + 0.057) \times (1 + 0.044) = 1.104$. This is implemented as an
intercept shift of $\ln(1.104)$, which preserves the ranking of
households.

\subsection{Validation}

Applied back to the 2023 source, the model predicts 15.61 percent
poverty against 15.50 percent observed, a close match. In L2Phl 2025 the
unadjusted prediction implies 16.96 percent poverty; after the GDP
adjustment, the preferred estimate is 13.20 percent.

\subsection{Output and link to weighting}

\begin{sloppypar}
The product is \rd{CAPI/S2S/target_hh_imputed.dta}, which holds a
per-capita imputed household income measure in the variable
\texttt{pcinc\_imp\_mean}. It is copied to
\rd{CATI/Analysis/HF/R00/target_hh_imputed.dta} and, in the CATI
analysis, merged onto the baseline weight file by household id to attach
imputed income for income-quintile breakdowns. It is an imputed-welfare
input, not a component of weight construction; Part~III documents how it
is merged.
\end{sloppypar}

\begin{notebox}{Caveat}
Updated 2025 regional poverty lines are not yet available, so the 2025
exercise uses the 2023 regional lines. If those lines rose with
inflation, the exercise may modestly understate 2025 poverty.
\end{notebox}

\begin{refbox}
\begin{sloppypar}
\rd{CAPI/S2S/S2S_PHL_note.docx} --- methodology (Sandra Segovia Juarez).
\rd{CAPI/S2S/s2s_diagnostics.xlsx} --- diagnostics.
\rd{CAPI/S2S/target_hh_imputed.dta} --- imputed output.
\end{sloppypar}
\end{refbox}

\section{Deliverables and Analytical Outputs}

The baseline findings are disseminated through a set of interactive HTML
products; this manual catalogs them and documents how they are rebuilt,
but does not reproduce their findings.

\begin{table}[htbp]
\centering
\begin{tabular}{@{}p{2.7cm}p{1.8cm}>{\raggedright\arraybackslash}p{10.4cm}@{}}
\toprule
Deliverable & Type & Location \\
\midrule
16-tab CAPI dashboard & dashboard & \rd{CAPI/Analysis/DB/html/l2phl_capi_dashboard.html} \\
12-chapter baseline storyline & storyline & \rd{CAPI/Analysis/SL/html/l2phl_baseline_story.html} \\
Wealth deep-dive chapter & storyline & \rd{CAPI/Analysis/SL/html/l2phl_wealth_chapter.html} \\
10 thematic deep dives & deep dives & \rd{CAPI/Analysis/DD/} \\
\bottomrule
\end{tabular}
\end{table}

\begin{sloppypar}
\noindent The individual deep-dive storylines live under each topic's
\texttt{html/} subfolder within \rd{CAPI/Analysis/DD/}. These products
are published on GitHub Pages at \url{https://avraltod.github.io/L2PHL/}.
\end{sloppypar}

\subsection{Production and Rebuild Pipeline}

Each deliverable is generated from a Stata do-file that writes results to
a machine-readable results file (\texttt{.md} or JSON), which a build
step then renders into a single self-contained HTML page (Chart.js via
CDN, Google Fonts). The storyline-rebuild and replication-audit workflows
regenerate a chapter from fresh Stata output, so the reported numbers
stay traceable back to the microdata.

\begin{refbox}
\rd{.claude/docs/storyline-update-workflow.md} --- storyline rebuild and
update workflow.
\end{refbox}
